\documentclass[11pt]{amsbook}
\usepackage[a4paper,margin=1in]{geometry}
\usepackage[T1]{fontenc}
\usepackage[utf8]{inputenc}
\usepackage{lmodern}
\usepackage{microtype}
\usepackage{amsmath,amssymb,amsthm,mathtools}
\usepackage{enumitem}
\usepackage{hyperref}
\usepackage{xcolor}
\hypersetup{colorlinks=true,linkcolor=blue!50!black,citecolor=blue!50!black,urlcolor=blue!50!black}

\newtheoremstyle{exstyle}%
  {1.2ex}% Space above
  {1.2ex}% Space below
  {\normalfont}% Body font
  {}% Indent amount
  {\bfseries}% Theorem head font
  {.}% Punctuation after theorem head
  {.75em}% Space after theorem head
  {}% Theorem head spec
\theoremstyle{exstyle}
\newtheorem{exercise}{Exercise}

\theoremstyle{plain}
\newtheorem*{remark}{Remark}

\newcommand{\F}{\mathbb{F}}
\newcommand{\Ccal}{\mathcal{C}}
\newcommand{\Dcal}{\mathcal{D}}
\newcommand{\Ecal}{\mathcal{E}}
\newcommand{\wt}{\operatorname{wt}}
\newcommand{\Span}{\operatorname{Span}}
\newcommand{\im}{\operatorname{im}}
\newcommand{\rk}{\operatorname{rk}}
\newcommand{\GL}{\operatorname{GL}}
\newcommand{\Aut}{\operatorname{Aut}}
\newcommand{\diag}{\operatorname{diag}}
\newcommand{\inner}[2]{\langle #1,#2\rangle}
\newcommand{\transpose}{\mathsf{T}}

\title{Exercises on Linear Codes\\\large A Book-Style Reorganization with Fifty Problems}
\author{}
\date{\today}

\begin{document}
\frontmatter
\maketitle
\tableofcontents

\chapter*{Preface}
These notes reorganize a homework-style collection of exercises on coding theory into a more standard mathematical format. The emphasis is on linear codes over finite fields, with a gradual progression from elementary constructions to duality, minimum distance, cosets, decoding, cyclic codes, and a first glimpse of algebraic-geometric codes.

The problems are intended to be read in order. Exercises~1--40 stay close to the classical core of an introductory course in coding theory. Exercises~41--50 extend the discussion toward more advanced topics.

\chapter{Exercises on Linear Codes}

\section{Notation and conventions}
Throughout, $\F_q$ denotes the finite field with $q$ elements. A linear code of length $n$ over $\F_q$ is a linear subspace $\Ccal \leq \F_q^n$. If $\dim_{\F_q}(\Ccal)=k$, then $\Ccal$ is called an \emph{$[n,k]_q$-code}. Its minimum Hamming distance is denoted by $d(\Ccal)$, and when this value is specified we write $[n,k,d]_q$.

If $G\in M_{k\times n}(\F_q)$ has row space $\Ccal$, then $G$ is a generator matrix for $\Ccal$. If $H\in M_{(n-k)\times n}(\F_q)$ satisfies
\[
\Ccal=\{x\in\F_q^n:Hx^{\transpose}=0\},
\]
then $H$ is a parity-check matrix for $\Ccal$. The dual code is
\[
\Ccal^{\perp}=\{y\in\F_q^n:\inner{x}{y}=0\text{ for all }x\in\Ccal\}.
\]
All vector spaces are over the ground field indicated in the statement.

\section{Basic constructions and systematic form}

\begin{exercise}[A single-erasure design problem]
Consider a communication channel over $\F_3$ that takes as input a vector of length $n=5$ and erases exactly one coordinate, without revealing in advance which coordinate will be erased. Construct two distinct linear $[5,4]_3$-codes that allow every message of length $4$ to be recovered from the received word.
\end{exercise}

\begin{exercise}[The repetition code]
Let $n\geq 1$, and let $\Ccal\subseteq \F_q^n$ be the repetition code
\[
\Ccal=\{(a,a,\dots,a):a\in\F_q\}.
\]
Determine its length, dimension, cardinality, and minimum distance. Write down a generator matrix and a parity-check matrix.
\end{exercise}

\begin{exercise}[The single parity-check code]
Let
\[
\Ccal=\Bigl\{(x_1,\dots,x_n)\in\F_q^n:\sum_{i=1}^n x_i=0\Bigr\}.
\]
Show that $\Ccal$ is a linear $[n,n-1]_q$-code. Determine its minimum distance, and exhibit both a generator matrix and a parity-check matrix.
\end{exercise}

\begin{exercise}[Systematic form over $\F_5$]
Let $\Ccal$ be the code generated over $\F_5$ by
\[
G=
\begin{bmatrix}
1&2&4&0&1\\
2&2&0&2&2\\
1&0&0&3&1
\end{bmatrix}.
\]
Use elementary row operations to rewrite $G$ in systematic form.
\end{exercise}

\begin{exercise}[Enumerating codewords over two fields]
Let $\Ccal$ be the code generated by
\[
G=
\begin{bmatrix}
1&0&1\\
0&1&1
\end{bmatrix}
\]
over $\F_q$.
\begin{enumerate}[label=\textup{(\arabic*)}]
\item For $q=2$, determine $|\Ccal|$ and list all codewords.
\item For $q=4$, determine $|\Ccal|$ and list all codewords, writing $\F_4=\{0,1,\omega,\omega^2\}$ with $\omega^2=\omega+1$.
\end{enumerate}
\end{exercise}

\begin{exercise}[Cardinality of a linear code]
Let $\Ccal\leq \F_q^n$ be a linear code of dimension $k$. Prove that $|\Ccal|=q^k$.
\end{exercise}

\begin{exercise}[Systematic encoding]
Suppose $G=[I_k\mid A]\in M_{k\times n}(\F_q)$ is a generator matrix in systematic form. Show that every information word $u\in\F_q^k$ produces a unique codeword $uG\in\Ccal$, and that the first $k$ coordinates of $uG$ are precisely the information symbols.
\end{exercise}

\begin{exercise}[A formal proof for systematic parity-check matrices]
Let
\[
G=[I_k\mid A]\in M_{k\times n}(\F_q).
\]
Prove that
\[
H=[-A^{\transpose}\mid I_{n-k}]
\]
is a parity-check matrix for the code generated by $G$.
\end{exercise}

\begin{exercise}[Length, dimension, and generator rank]
Let $G\in M_{k\times n}(\F_q)$ be a matrix of rank $r$. Show that the row space of $G$ is a linear code of length $n$, dimension $r$, and cardinality $q^r$.
\end{exercise}

\begin{exercise}[Puncturing and shortening]
Let $\Ccal\leq\F_q^n$ be a linear code, and fix a coordinate position $i$. Define the punctured code $\Ccal^{(i)}$ by deleting the $i$th coordinate from every codeword. Show that $\Ccal^{(i)}$ is a linear code of length $n-1$, and prove that
\[
\dim \Ccal -1 \leq \dim \Ccal^{(i)} \leq \dim \Ccal.
\]
\end{exercise}

\section{Duality and self-duality}

\begin{exercise}[A parity-check matrix in systematic form]
Consider the code generated by
\[
G=
\begin{bmatrix}
1&0&0&0&0&1&1&1\\
0&1&0&0&1&0&1&1\\
0&0&1&0&1&1&0&1\\
0&0&0&1&1&1&1&0
\end{bmatrix}.
\]
Compute a parity-check matrix in systematic form.
\end{exercise}

\begin{exercise}[Self-orthogonal or self-dual?]
Let $\Ccal\leq \F_2^8$ be the code generated by
\[
G=
\begin{bmatrix}
1&0&0&0&0&1&1&1\\
0&1&0&0&1&0&1&1\\
0&0&1&0&1&1&0&1\\
0&0&0&1&1&1&1&0
\end{bmatrix}.
\]
Determine whether $\Ccal$ is self-orthogonal and whether it is self-dual. Decide whether the answer changes if one replaces $\F_2$ by $\F_p$, where $p$ is an odd prime.
\end{exercise}

\begin{exercise}[A necessary condition for self-duality]
Show that if a linear code $\Ccal\leq \F_q^n$ is self-dual, then $n$ is even and $\dim \Ccal=n/2$.
\end{exercise}

\begin{exercise}[The tetracode from a parity-check matrix]
Let $\Ccal_1$ be the code over $\F_3$ with parity-check matrix
\[
H_1=
\begin{bmatrix}
2&2&1&0\\
0&1&1&2
\end{bmatrix}.
\]
Determine whether $\Ccal_1$ is self-orthogonal and whether it is self-dual.
\end{exercise}

\begin{exercise}[A code over $\F_4$]
Let $\omega\in \F_4$ satisfy $\omega^2=\omega+1$, and let $\Ccal_2$ be the code generated by
\[
G_2=
\begin{bmatrix}
1&0&0&1&\omega&\omega\\
0&1&0&\omega&1&\omega\\
0&0&1&\omega&\omega&1
\end{bmatrix}.
\]
Determine whether $\Ccal_2$ is self-orthogonal and whether it is self-dual.
\end{exercise}

\begin{exercise}[Weight and self-inner product over $\F_3$]
For $x\in \F_3^n$, prove that
\[
\wt(x)\equiv \inner{x}{x}\pmod 3.
\]
\end{exercise}

\begin{exercise}[Doubly-even type argument]
Let $\Ccal\leq\F_2^n$ be a self-orthogonal code having a generator matrix whose rows all have Hamming weight divisible by $4$. Prove that every codeword of $\Ccal$ has Hamming weight divisible by $4$.
\end{exercise}

\begin{exercise}[Dual of a sum]
For linear codes $\Ccal_1,\Ccal_2\leq\F_q^n$, prove that
\[
(\Ccal_1+\Ccal_2)^{\perp}=\Ccal_1^{\perp}\cap \Ccal_2^{\perp}.
\]
\end{exercise}

\begin{exercise}[Dual of an intersection]
For linear codes $\Ccal_1,\Ccal_2\leq\F_q^n$, prove that
\[
(\Ccal_1\cap \Ccal_2)^{\perp}=\Ccal_1^{\perp}+\Ccal_2^{\perp}.
\]
\end{exercise}

\begin{exercise}[Dimension formula]
Let $\Ccal\leq\F_q^n$ be a linear code. Prove that
\[
\dim \Ccal+\dim \Ccal^{\perp}=n.
\]
Deduce that $\Ccal\subseteq \Ccal^{\perp}$ if and only if $\dim \Ccal\leq n/2$.
\end{exercise}

\section{Minimum distance and parity-check geometry}

\begin{exercise}[Distance of the single parity-check code]
Using only the parity-check matrix $H=[1\;1\;\cdots\;1]$, determine the minimum distance of the $[n,n-1]_q$ single parity-check code.
\end{exercise}

\begin{exercise}[Distance of the tetracode]
The tetracode over $\F_3$ is
\[
\Ccal=\{(x_1,x_2,x_1+x_2,x_1-x_2):x_1,x_2\in\F_3\}.
\]
Using a parity-check matrix for this code, determine its minimum distance.
\end{exercise}

\begin{exercise}[A binary $(7,4)$ code]
Let $\Ccal\leq \F_2^7$ be generated by
\[
G=
\begin{bmatrix}
1&0&0&0&1&0&1\\
0&1&0&0&1&1&0\\
0&0&1&0&0&1&1\\
0&0&0&1&1&1&1
\end{bmatrix}.
\]
Determine the minimum distance of $\Ccal$.
\end{exercise}

\begin{exercise}[A binary $(6,3)$ code]
Let $\Ccal\leq\F_2^6$ be generated by
\[
G=
\begin{bmatrix}
1&1&0&1&1&0\\
0&1&0&1&0&1\\
1&1&1&0&0&0
\end{bmatrix}.
\]
Determine the minimum distance of $\Ccal$, and decide whether $\Ccal$ contains a codeword of weight $4$.
\end{exercise}

\begin{exercise}[Distance from a parity-check matrix]
Let $\Ccal\leq\F_q^n$ have parity-check matrix $H$. Prove that $d(\Ccal)$ is the smallest number of columns of $H$ that are linearly dependent.
\end{exercise}

\begin{exercise}[Equivalent definitions of minimum distance]
Let $\Ccal\leq \F_q^n$ be a nonzero linear code. Prove that
\[
d(\Ccal)=\min\{\wt(c):0\neq c\in\Ccal\}.
\]
\end{exercise}

\begin{exercise}[Error-correcting capability]
Let $\Ccal$ be a code of minimum distance $d$. Prove that every Hamming ball of radius
\[
t=\Bigl\lfloor\frac{d-1}{2}\Bigr\rfloor
\]
centered at a codeword is disjoint from every other such ball.
\end{exercise}

\begin{exercise}[Perfect repetition codes]
Let $\Ccal\subseteq\F_2^n$ be the binary repetition code of length $n$. Determine for which values of $n$ this code is perfect.
\end{exercise}

\begin{exercise}[Puncturing a binary $(24,12,8)$ code]
Assume that a binary $[24,12,8]$ code exists. Does puncturing it in one coordinate yield a perfect code? Justify your answer using the sphere-packing bound.
\end{exercise}

\begin{exercise}[Extending a binary $(5,3,3)$ code]
Suppose a binary $[5,3,3]$ code is extended by one coordinate. Determine whether the resulting code can be perfect. Give a geometric interpretation of perfect codes in terms of Hamming spheres.
\end{exercise}

\section{Cosets, syndromes, and decoding}

\begin{exercise}[Disjoint Hamming spheres in the tetracode]
Let $\Ccal$ be the tetracode over $\F_3$:
\[
\Ccal=\{(x_1,x_2,x_1+x_2,x_1-x_2):x_1,x_2\in\F_3\}.
\]
Construct the Hamming spheres of radius $1$ around the codewords of $\Ccal$ and prove that they are pairwise disjoint.
\end{exercise}

\begin{exercise}[A non-linear subset inside a linear code]
Let $\Ccal\leq\F_2^6$ be generated by
\[
G=
\begin{bmatrix}
1&1&0&0&0&0\\
0&1&1&0&0&0\\
1&1&1&1&1&1
\end{bmatrix}.
\]
Decide whether the set of codewords of $\Ccal$ whose weights are divisible by $4$ forms a linear code.
\end{exercise}

\begin{exercise}[A coset-weight construction]
Find a nonzero binary linear $[4,k]$ code $\Ccal$ such that every coset of $\Ccal$ contains a unique vector of minimum weight, while some coset has weight strictly greater than
\[
\Bigl\lfloor\frac{d(\Ccal)-1}{2}\Bigr\rfloor.
\]
\end{exercise}

\begin{exercise}[Syndrome decoding for the tetracode]
Let $\Ccal$ be the tetracode over $\F_3$. Decode the received vectors
\[
(1,1,1,1),\qquad (1,-1,0,-1)
\]
by means of a syndrome table.
\end{exercise}

\begin{exercise}[Cosets of a binary $(7,4)$ code]
For the code of Exercise~23:
\begin{enumerate}[label=\textup{(\arabic*)}]
\item determine the number of cosets,
\item construct a complete table of syndromes and coset leaders,
\item decode the received vectors $(1,1,0,1,0,1,0)$ and $(0,0,1,0,0,0,0)$.
\end{enumerate}
\end{exercise}

\begin{exercise}[A systematic analysis of a binary $(6,3)$ code]
For the code of Exercise~24:
\begin{enumerate}[label=\textup{(\arabic*)}]
\item rewrite the generator matrix in systematic form,
\item compute a parity-check matrix,
\item decide whether the code is self-orthogonal or self-dual,
\item determine the number of cosets.
\end{enumerate}
\end{exercise}

\begin{exercise}[Erasure decoding]
For the binary $[6,3]$ code of Exercise~24, suppose the decoder receives
\[
(1,0,0,0,\ast,1),
\]
where $\ast$ denotes an erasure. Determine the transmitted codeword.
\end{exercise}

\begin{exercise}[A coset of weight $2$]
For the binary $[6,3]$ code of Exercise~24, identify a coset of weight $2$ and determine all of its coset leaders.
\end{exercise}

\begin{exercise}[When is syndrome decoding unambiguous?]
Let $\Ccal\leq\F_q^n$ have minimum distance $d$. Prove that every coset has at most one leader of weight at most $\lfloor (d-1)/2\rfloor$. Explain why this is the theoretical basis of unique decoding up to the correction radius.
\end{exercise}

\begin{exercise}[Syndromes and cosets]
Let $H$ be a parity-check matrix of a linear code $\Ccal\leq\F_q^n$. Prove that two words $x,y\in\F_q^n$ have the same syndrome if and only if they lie in the same coset of $\Ccal$.
\end{exercise}

\section{Further exercises and extensions}

\begin{exercise}[Monomial equivalence and duals]
Give an example of two codes $\Ccal_1$ and $\Ccal_2$ over the same field, together with a monomial matrix $M$, such that
\[
\Ccal_1M=\Ccal_2
\qquad\text{but}\qquad
\Ccal_1^{\perp}M\neq \Ccal_2^{\perp}.
\]
\end{exercise}

\begin{exercise}[A binary weight identity]
If $x,y\in\F_2^n$, let $x\cap y$ denote the vector whose support is the intersection of the supports of $x$ and $y$. Prove that
\[
\wt(x+y)=\wt(x)+\wt(y)-2\wt(x\cap y).
\]
\end{exercise}

\begin{exercise}[The ternary Golay code $G_{12}$]
Let $G_{12}$ be the ternary code generated by $[I_6\mid A]$, where
\[
A=
\begin{bmatrix}
0&1&1&1&1&1\\
1&0&1&2&2&1\\
1&1&0&1&2&2\\
1&2&1&0&1&2\\
1&2&2&1&0&1\\
1&1&2&2&1&0
\end{bmatrix}.
\]
Show that $G_{12}$ is a self-dual $[12,6,6]_3$-code. Then explain why puncturing an appropriate coordinate produces an $[11,6,5]_3$ code.
\end{exercise}

\begin{exercise}[Puncturing, extending, and even-like codewords]
Let $\Ccal\leq\F_q^n$ be a linear code, and let $\Ccal_1$ be obtained by puncturing the rightmost coordinate and then extending the punctured code by appending the negative sum of the remaining coordinates.
\begin{enumerate}[label=\textup{(\arabic*)}]
\item Prove that $\Ccal=\Ccal_1$ if and only if every codeword of $\Ccal$ has coordinate sum $0$.
\item Prove that if $\Ccal$ is self-orthogonal and contains the all-one vector, then $\Ccal=\Ccal_1$.
\end{enumerate}
\end{exercise}

\begin{exercise}[A first Reed--Solomon exercise]
Let $\alpha_1,\dots,\alpha_n\in\F_q$ be distinct, and let $1\leq k\leq n$. Define
\[
\operatorname{RS}_{k}(\alpha_1,\dots,\alpha_n)
 =\{(f(\alpha_1),\dots,f(\alpha_n)): f\in \F_q[x],\ \deg f<k\}.
\]
Show that this is a linear $[n,k]_q$-code.
\end{exercise}

\begin{exercise}[Reed--Solomon codes are MDS]
With the notation of Exercise~45, prove that
\[
d\bigl(\operatorname{RS}_k(\alpha_1,\dots,\alpha_n)\bigr)=n-k+1.
\]
\end{exercise}

\begin{exercise}[Cyclic codes as ideals]
Assume $\gcd(n,q)=1$. Let $R=\F_q[x]/\langle x^n-1\rangle$, and let $\sigma$ denote the cyclic shift on $\F_q^n$:
\[
\sigma(c_0,c_1,\dots,c_{n-1})=(c_{n-1},c_0,\dots,c_{n-2}).
\]
Show that the linear codes $\Ccal\leq\F_q^n$ that are invariant under $\sigma$ correspond exactly to ideals of the quotient ring $R$.
\end{exercise}

\begin{exercise}[Invariant subspaces and the quotient-ring model]
Make explicit the correspondence of Exercise~47 by sending
\[
(c_0,c_1,\dots,c_{n-1})\longmapsto c_0+c_1x+\cdots+c_{n-1}x^{n-1}\in R.
\]
Prove that invariance under cyclic permutation is equivalent to closure under multiplication by $x$ modulo $x^n-1$.
\end{exercise}

\begin{exercise}[A binary Goppa code identity]
Let $L=\{\alpha_1,\dots,\alpha_n\}\subseteq \F_{2^m}$, and let $g(x)\in \F_{2^m}[x]$ be squarefree with $g(\alpha_i)\neq 0$ for all $i$. The classical binary Goppa code $\Gamma(L,g)$ may be described as
\[
\Gamma(L,g)=\Bigl\{c=(c_i)\in\F_2^n:\sum_{i=1}^n \frac{c_i}{x-\alpha_i}\equiv 0\pmod{g(x)}\Bigr\}.
\]
Show that
\[
\Gamma(L,g)=\Gamma(L,g^2).
\]
\end{exercise}

\begin{exercise}[A first algebraic-geometric coding exercise]
Let $X$ be a smooth projective curve of genus $g$ over $\F_q$, let
\[
D=P_1+\cdots+P_n
\]
be a sum of distinct rational points on $X$, and let $G$ be a divisor whose support is disjoint from that of $D$. Define the evaluation code
\[
C_L(D,G)=\{(f(P_1),\dots,f(P_n)):f\in L(G)\}.
\]
\begin{enumerate}[label=\textup{(\arabic*)}]
\item Using the Riemann--Roch theorem, prove that if $\deg G<n$, then
\[
\dim C_L(D,G)\geq \deg G+1-g.
\]
If in addition $2g-2<\deg G<n$, prove that equality holds.
\item Prove that every nonzero codeword of $C_L(D,G)$ has at most $\deg G$ zero coordinates. Deduce that
\[
d\bigl(C_L(D,G)\bigr)\geq n-\deg G.
\]
\end{enumerate}
\end{exercise}

\backmatter
\chapter*{Concluding remark}
This set is intended as a bridge between computational fluency and conceptual structure. The first forty problems may be approached directly with row reduction, parity-check matrices, and elementary counting arguments. The final ten problems indicate how the same language continues into Reed--Solomon, cyclic, Goppa, and algebraic-geometric codes.

\end{document}
